\sscp{Different sound correspondences in language/dialect clusters}{02-04-07}
Different varieties of larger “languages” or clusters of closely
related languages and dialects often have different historical sound correspondences.
This is true, for example,
for the languages of Rote Island
and the \tsc{Meto} languages of western Timor.
Many of these systematic differences have been known for
some time \citep{Jonker1908,Jonker1915}.
These systematic distinctions are very important for unravelling
the prehistory of the region \citep{Edwards2018a,Edwards2018b,Edwards2021}.
Yet comparative sources that simply say “Rote” (or “Roti” [sic])
or “Uab \ili{Meto}” (or “Dawan” or “Atoni’”),
and fail to specify which variety,
may mask important differences,
rather than track systematic patterns rigorously.
For example, PMP *\tbr{z}a\tbb{l}an ‘path, way’ >
Roi{\Q}is \ili{Amarasi} \it{\tbr{r}a\tbb{n}an},
%\ili{Kotos} \ili{Amarasi} \it{\tbr{r}a\tbb{n}an},
\ili{Amanuban} \it{\tbr{l}a\tbb{n}an},
and Amfo{\Q}an \it{\tbr{l}a\tbb{l}an} reflecting different correspondences
of PMP *z and *l.
In the past these have often been treated as a single language
in spite of cumulatively significant differences in the sound
correspondences, lexicon, morphophonemic processes, and grammatical functors.
The \tsc{\ili{Ambon}-Seram} language \ili{Alune} spoken on Seram also has
considerable variation in its many varieties
and in their historical sound correspondences.